% Proofs for Section 14: Abstract environments and selective sensing (key env)
% src: proofs of H5 (H5-L1 to H5-L5, H5-T1 to H5-T3, H5-P1 to H5-P7, the lemma of H5-D3), as repaired by critic V11 and Section 2
\section{Proofs for \texorpdfstring{Section~\ref{sec:env}}{abstract environments}}\label{app:env}

The proofs follow the order of \cref{sec:env}. In the setting of \cref{def:env-sensor} put $\tilde{\mathsf X}=\Sigma_{\mathsf X}^{-1/2}\mathsf X$ and $\tilde{\mathsf Y}=\Sigma_{\mathsf Y}^{-1/2}\mathsf Y$, so that $\Cov(\tilde{\mathsf X})=I$, $\Cov(\tilde{\mathsf X},\tilde{\mathsf Y})=\mathsf C$ and $s=\tilde H\tilde{\mathsf X}+v$. For a symmetric matrix, $\lambda^\downarrow_{\mathsf j}$ and $\lambda^\uparrow_{\mathsf j}$ are its eigenvalues in decreasing and in increasing order. An orthogonal change of a sensor's channels, $\tilde H\mapsto\mathsf O\tilde H$, leaves $\mathsf G_s$ unchanged, and after a singular value decomposition the channels beyond the rank of $\tilde H$ are pure noise; we may therefore take $\kappa\le d_{\mathsf X}$ whenever convenient.

\begin{lemma}[Eigenvalue inequalities]\label{lem:env-eigen}
\begin{enumerate}
\item[(a)] (Interlacing.) For symmetric $\mathsf S\in\R^{d\times d}$ and $\mathsf U\in\R^{d\times\mathsf m}$ with $\mathsf U^\top\mathsf U=I_{\mathsf m}$, $\lambda^\downarrow_{\mathsf j}(\mathsf U^\top\mathsf S\mathsf U)\le\lambda^\downarrow_{\mathsf j}(\mathsf S)$ and $\lambda^\uparrow_{\mathsf j}(\mathsf U^\top\mathsf S\mathsf U)\ge\lambda^\uparrow_{\mathsf j}(\mathsf S)$ for $\mathsf j\le\mathsf m$.
\item[(b)] (Ky Fan.) $\max\{\tr(\mathsf U^\top\mathsf S\mathsf U):\mathsf U^\top\mathsf U=I_{\mathsf m}\}=\sum_{\mathsf j\le\mathsf m}\lambda^\downarrow_{\mathsf j}(\mathsf S)$, and if $\lambda^\downarrow_{\mathsf m}(\mathsf S)>\lambda^\downarrow_{\mathsf m+1}(\mathsf S)$ the maximizers are the $\mathsf U$ whose columns span the top-$\mathsf m$ eigenspace.
\item[(c)] (Concave majorization.) Let $x,y\in\R^d$ be non-increasing, with $\sum_{\mathsf j\le\mathsf m}x_{\mathsf j}\le\sum_{\mathsf j\le\mathsf m}y_{\mathsf j}$ for $\mathsf m<d$ and $\sum_{\mathsf j\le d}x_{\mathsf j}=\sum_{\mathsf j\le d}y_{\mathsf j}$, and let $\varphi$ be concave and differentiable on an open interval that contains all entries. Then $\sum_{\mathsf j}\varphi(x_{\mathsf j})\ge\sum_{\mathsf j}\varphi(y_{\mathsf j})$.
\item[(d)] (Determinant of a sum.) For $\mathsf S_1,\mathsf S_2\succeq0$ in $\R^{d\times d}$, $\det(\mathsf S_1+\mathsf S_2)\ge\prod_{\mathsf j}\bigl(\lambda^\downarrow_{\mathsf j}(\mathsf S_1)+\lambda^\downarrow_{\mathsf j}(\mathsf S_2)\bigr)$.
\item[(e)] (Trace bound.) For $\mathsf G\succeq0$ of rank at most $\kappa$ with eigenvalues $\mathsf g_1\ge\dots\ge\mathsf g_\kappa$ and symmetric $\mathsf J$, $\tr(\mathsf G\mathsf J)\le\sum_{\mathsf j\le\kappa}\mathsf g_{\mathsf j}\lambda^\downarrow_{\mathsf j}(\mathsf J)$, with equality when the eigenvector of each positive $\mathsf g_{\mathsf j}$ is an eigenvector of $\lambda^\downarrow_{\mathsf j}(\mathsf J)$.
\item[(f)] (Determinant bound for sensors.) For $\tilde H\in\R^{\kappa\times d}$, $\kappa\le d$, with singular values $\mathsf d_1\ge\dots\ge\mathsf d_\kappa\ge0$, and $\mathsf S\succeq0$ in $\R^{d\times d}$, $\det(I_\kappa+\tilde H\mathsf S\tilde H^\top)\ge\prod_{\mathsf j\le\kappa}\bigl(1+\mathsf d_{\mathsf j}^2\lambda^\uparrow_{\mathsf j}(\mathsf S)\bigr)$, with equality when $\tilde H=\mathsf O\diag(\mathsf d)[\mathsf f_1,\dots,\mathsf f_\kappa]^\top$, $\mathsf O$ orthogonal and $\mathsf f_{\mathsf j}$ orthonormal eigenvectors of $\lambda^\uparrow_{\mathsf j}(\mathsf S)$.
\end{enumerate}
\end{lemma}
% src: H5-L2 (the attribution of (d) dropped; critic V11)

\begin{proof}
(a) By the Courant--Fischer theorem \citep{horn2012matrix}, $\lambda^\downarrow_{\mathsf j}(\mathsf U^\top\mathsf S\mathsf U)$ is the maximum over $\mathsf j$-dimensional subspaces $\mathcal V\subseteq\R^{\mathsf m}$ of $\min\{x^\top\mathsf U^\top\mathsf S\mathsf Ux:x\in\mathcal V,\lvert x\rvert=1\}$; since $\mathsf U\mathcal V$ is a $\mathsf j$-dimensional subspace of $\R^d$ and $\lvert\mathsf Ux\rvert=\lvert x\rvert$, this is at most $\lambda^\downarrow_{\mathsf j}(\mathsf S)$. Applied to $-\mathsf S$, the same argument gives the second inequality.

(b) Write $\mathsf S=\sum_{\mathsf j}\mathsf s_{\mathsf j}\mathsf b_{\mathsf j}\mathsf b_{\mathsf j}^\top$ with $\mathsf s_{\mathsf j}=\lambda^\downarrow_{\mathsf j}(\mathsf S)$ and orthonormal $\mathsf b_{\mathsf j}$, and put $\mathsf y_{\mathsf j}=\lvert\mathsf U^\top\mathsf b_{\mathsf j}\rvert^2\in[0,1]$, so that $\sum_{\mathsf j}\mathsf y_{\mathsf j}=\tr(\mathsf U\mathsf U^\top)=\mathsf m$ and $\tr(\mathsf U^\top\mathsf S\mathsf U)=\sum_{\mathsf j}\mathsf s_{\mathsf j}\mathsf y_{\mathsf j}$. Since $\sum_{\mathsf j\le\mathsf m}(1-\mathsf y_{\mathsf j})=\sum_{\mathsf j>\mathsf m}\mathsf y_{\mathsf j}$,
\[
\sum_{\mathsf j}\mathsf s_{\mathsf j}\mathsf y_{\mathsf j}-\sum_{\mathsf j\le\mathsf m}\mathsf s_{\mathsf j}=-\sum_{\mathsf j\le\mathsf m}\mathsf s_{\mathsf j}(1-\mathsf y_{\mathsf j})+\sum_{\mathsf j>\mathsf m}\mathsf s_{\mathsf j}\mathsf y_{\mathsf j}\le(\mathsf s_{\mathsf m+1}-\mathsf s_{\mathsf m})\sum_{\mathsf j>\mathsf m}\mathsf y_{\mathsf j}\le0.
\]
Equality holds for $\mathsf U=[\mathsf b_1,\dots,\mathsf b_{\mathsf m}]$, and under the gap it forces $\sum_{\mathsf j>\mathsf m}\mathsf y_{\mathsf j}=0$, so the columns of $\mathsf U$ span the top-$\mathsf m$ eigenspace.

(c) Put $\mathsf a_{\mathsf j}=\int_0^1\varphi'\bigl((1-\vartheta)y_{\mathsf j}+\vartheta x_{\mathsf j}\bigr)\mathrm d\vartheta$, so that $\varphi(x_{\mathsf j})-\varphi(y_{\mathsf j})=\mathsf a_{\mathsf j}(x_{\mathsf j}-y_{\mathsf j})$. For each $\vartheta$ the point $(1-\vartheta)y_{\mathsf j}+\vartheta x_{\mathsf j}$ is non-increasing in $\mathsf j$ and $\varphi'$ is non-increasing, so $\mathsf a_1\le\mathsf a_2\le\dots\le\mathsf a_d$. With $\mathsf t_{\mathsf m}=\sum_{\mathsf j\le\mathsf m}(x_{\mathsf j}-y_{\mathsf j})\le0$ and $\mathsf t_d=0$, Abel summation gives $\sum_{\mathsf j}\mathsf a_{\mathsf j}(x_{\mathsf j}-y_{\mathsf j})=\sum_{\mathsf m<d}(\mathsf a_{\mathsf m}-\mathsf a_{\mathsf m+1})\mathsf t_{\mathsf m}\ge0$.

(d) Let $x=\lambda^\downarrow(\mathsf S_1+\mathsf S_2)$ and $y_{\mathsf j}=\lambda^\downarrow_{\mathsf j}(\mathsf S_1)+\lambda^\downarrow_{\mathsf j}(\mathsf S_2)$, both non-increasing. By (b), $\sum_{\mathsf j\le\mathsf m}x_{\mathsf j}=\max_{\mathsf U}\tr(\mathsf U^\top(\mathsf S_1+\mathsf S_2)\mathsf U)\le\sum_{\mathsf j\le\mathsf m}y_{\mathsf j}$, with equality at $\mathsf m=d$, where both sides are traces. If $y_d=0$ the right side of (d) vanishes. Otherwise $x_d=\tr(\mathsf S_1+\mathsf S_2)-\sum_{\mathsf j<d}x_{\mathsf j}\ge y_d>0$, all entries are positive, and (c) with $\varphi=\ln$ gives $\sum_{\mathsf j}\ln x_{\mathsf j}\ge\sum_{\mathsf j}\ln y_{\mathsf j}$.

(e) Write $\mathsf G=\sum_{\mathsf j\le\kappa}\mathsf g_{\mathsf j}\mathsf f_{\mathsf j}\mathsf f_{\mathsf j}^\top$ with orthonormal $\mathsf f_{\mathsf j}$, and put $\mathsf g_{\kappa+1}=0$. Abel summation and (b) give
\[
\tr(\mathsf G\mathsf J)=\sum_{\mathsf m\le\kappa}(\mathsf g_{\mathsf m}-\mathsf g_{\mathsf m+1})\sum_{\mathsf j\le\mathsf m}\mathsf f_{\mathsf j}^\top\mathsf J\mathsf f_{\mathsf j}\le\sum_{\mathsf m\le\kappa}(\mathsf g_{\mathsf m}-\mathsf g_{\mathsf m+1})\sum_{\mathsf j\le\mathsf m}\lambda^\downarrow_{\mathsf j}(\mathsf J)=\sum_{\mathsf j\le\kappa}\mathsf g_{\mathsf j}\lambda^\downarrow_{\mathsf j}(\mathsf J),
\]
with equality in the stated case.

(f) Factors with $\mathsf d_{\mathsf j}=0$ equal one on both sides, so let $\kappa'$ be the number of positive singular values and write $\tilde H=\mathsf O'\mathsf H_{\mathrm d}\mathsf U^\top$ with $\mathsf O'$ and $\mathsf U$ having $\kappa'$ orthonormal columns and $\mathsf H_{\mathrm d}=\diag(\mathsf d_1,\dots,\mathsf d_{\kappa'})$. By Sylvester's identity, $\det(I_\kappa+\tilde H\mathsf S\tilde H^\top)=\det(I+\mathsf H_{\mathrm d}\mathsf N\mathsf H_{\mathrm d})=\det(\mathsf H_{\mathrm d})^2\det(\mathsf H_{\mathrm d}^{-2}+\mathsf N)$ with $\mathsf N=\mathsf U^\top\mathsf S\mathsf U$. By (d), $\det(\mathsf H_{\mathrm d}^{-2}+\mathsf N)\ge\prod_{\mathsf j}\bigl(\lambda^\downarrow_{\mathsf j}(\mathsf H_{\mathrm d}^{-2})+\lambda^\downarrow_{\mathsf j}(\mathsf N)\bigr)$; since $\lambda^\downarrow_{\mathsf j}(\mathsf H_{\mathrm d}^{-2})=\mathsf d_{\kappa'+1-\mathsf j}^{-2}$ and $\lambda^\downarrow_{\mathsf j}(\mathsf N)=\lambda^\uparrow_{\kappa'+1-\mathsf j}(\mathsf N)$, the product is $\prod_{\mathsf j}(\mathsf d_{\mathsf j}^{-2}+\lambda^\uparrow_{\mathsf j}(\mathsf N))$, and multiplying by $\det(\mathsf H_{\mathrm d})^2$ gives $\prod_{\mathsf j}(1+\mathsf d_{\mathsf j}^2\lambda^\uparrow_{\mathsf j}(\mathsf N))$. By (a), $\lambda^\uparrow_{\mathsf j}(\mathsf N)\ge\lambda^\uparrow_{\mathsf j}(\mathsf S)$. In the equality case $\mathsf N=\diag(\lambda^\uparrow_{\mathsf j}(\mathsf S))$ and the determinant is the product.
\end{proof}

\begin{proof}[Proof of \cref{lem:env-information}]
Gaussian conditioning gives $\E[\tilde{\mathsf X}\mid s]=\tilde H^\top(I+\tilde H\tilde H^\top)^{-1}s$, whose covariance is $\mathsf G_s$, so $\Cov(\tilde{\mathsf X}\mid s)=I-\mathsf G_s$. With the singular value decomposition $\tilde H=\mathsf O\diag(\mathsf d)\mathsf U^\top$, $\mathsf G_s=\mathsf U\diag(\mathsf d_{\mathsf j}^2/(1+\mathsf d_{\mathsf j}^2))\mathsf U^\top$, which gives the captured fractions, and (i) is $I(s;\tilde{\mathsf X})=h(\tilde{\mathsf X})-h(\tilde{\mathsf X}\mid s)$. (ii) Since $v$ is independent of $(\mathsf X,\mathsf Y)$, $\E[\tilde{\mathsf Y}\mid\tilde{\mathsf X},s]=\mathsf C^\top\tilde{\mathsf X}$, hence $\E[\tilde{\mathsf Y}\mid s]=\mathsf C^\top\E[\tilde{\mathsf X}\mid s]$, whose covariance is $\mathsf C^\top\mathsf G_s\mathsf C$; Gaussian conditional covariances do not depend on the data, so $\Cov(\tilde{\mathsf Y}\mid s)=I-\mathsf C^\top\mathsf G_s\mathsf C$, and $I(s;\mathsf Y)=h(\tilde{\mathsf Y})-h(\tilde{\mathsf Y}\mid s)$. Since $s$ and $\mathsf Y$ are conditionally independent given $\mathsf X$, the chain rule also gives $I(s;\mathsf Y)=I(s;\mathsf X)-I(s;\mathsf X\mid\mathsf Y)$, and with $\Cov(\tilde{\mathsf X}\mid\mathsf Y)=I-\mathsf C\mathsf C^\top$,
\begin{equation}\label{eq:env-second-form}
I(s;\mathsf Y)=\tfrac12\ln\det\bigl(I+\tilde H\tilde H^\top\bigr)-\tfrac12\ln\det\bigl(I+\tilde H(I-\mathsf C\mathsf C^\top)\tilde H^\top\bigr).
\end{equation}
(iii) $\Cov(\mathsf Y\mid s)=\Sigma_{\mathsf Y}^{1/2}(I-\mathsf C^\top\mathsf G_s\mathsf C)\Sigma_{\mathsf Y}^{1/2}$ and $\Sigma_{\mathsf Y}^{1/2}\mathsf C^\top=\Sigma_{\mathsf Y\mathsf X}\Sigma_{\mathsf X}^{-1/2}$. (iv) Captured covariances have the stated form by the decomposition above. Conversely, for $\mathsf G=\sum_{\mathsf j\le\kappa}\mathsf g_{\mathsf j}\mathsf f_{\mathsf j}\mathsf f_{\mathsf j}^\top$ with $\mathsf g_{\mathsf j}\in[0,1)$ and orthonormal $\mathsf f_{\mathsf j}$, the whitened sensor with rows $(\mathsf g_{\mathsf j}/(1-\mathsf g_{\mathsf j}))^{1/2}\mathsf f_{\mathsf j}^\top$ has captured covariance $\mathsf G$.
\end{proof}

\begin{proof}[Proof of \cref{thm:env-canonical}]
(a) Apply \cref{lem:env-eigen}(f) to the second determinant in \eqref{eq:env-second-form}, with $\mathsf S=I-\mathsf C\mathsf C^\top$, whose eigenvalues in increasing order are $1-\sigma_1^2\le1-\sigma_2^2\le\dots$. Since $\det(I+\tilde H\tilde H^\top)=\prod_{\mathsf j}(1+\mathsf d_{\mathsf j}^2)$,
\[
I(s;\mathsf Y)\le\sum_{\mathsf j\le\kappa}\tfrac12\ln\frac{1+\mathsf d_{\mathsf j}^2}{1+\mathsf d_{\mathsf j}^2(1-\sigma_{\mathsf j}^2)}=-\tfrac12\sum_{\mathsf j\le\kappa}\ln(1-\mathsf g_{\mathsf j}\sigma_{\mathsf j}^2),
\]
because $\mathsf g_{\mathsf j}=\mathsf d_{\mathsf j}^2/(1+\mathsf d_{\mathsf j}^2)$. The equality case of \cref{lem:env-eigen}(f) is the sensor whose $\mathsf j$-th row of $\tilde H$ is $\mathsf d_{\mathsf j}\mathsf e_{\mathsf j}^\top$.

(b) Let $\mathsf c_{\mathsf j}=-\frac12\ln(1-\mathsf g_{\mathsf j})\ge0$ be the capacity spent on channel $\mathsf j$. By (a) and \cref{lem:env-information}(i), every sensor satisfies $I(s;\mathsf Y)\le\sum_{\mathsf j\le\kappa}\mathcal Y_{\mathsf j}(\mathsf c_{\mathsf j})$ with $\sum_{\mathsf j}\mathsf c_{\mathsf j}=I(s;\mathsf X)$ and $\mathcal Y_{\mathsf j}(x)=-\frac12\ln(1-\sigma_{\mathsf j}^2(1-e^{-2x}))$, and the canonical sensor attains this bound. The optimum is therefore the maximum of $\sum_{\mathsf j\le\kappa}\mathcal Y_{\mathsf j}(\mathsf c_{\mathsf j})$ over $\mathsf c_{\mathsf j}\ge0$ with $\sum_{\mathsf j}\mathsf c_{\mathsf j}\le I_{\max}$, provided the maximizer is non-increasing in $\mathsf j$, as the one found below is, so that its pairing with the ordered correlations in (a) costs nothing. For $0<\sigma_{\mathsf j}<1$ the derivative is
\[
\mathcal Y_{\mathsf j}'(x)=\frac{\sigma_{\mathsf j}^2}{(1-\sigma_{\mathsf j}^2)e^{2x}+\sigma_{\mathsf j}^2},
\]
which is positive and decreasing, so each $\mathcal Y_{\mathsf j}$ is increasing and concave ($\mathcal Y_{\mathsf j}\equiv0$ when $\sigma_{\mathsf j}=0$). For a separable concave objective on this polytope the Karush--Kuhn--Tucker conditions are necessary and sufficient: there is $\mathsf v>0$ with $\mathcal Y_{\mathsf j}'(\mathsf c_{\mathsf j})=\mathsf v$ when $\mathsf c_{\mathsf j}>0$ and $\mathcal Y_{\mathsf j}'(0)=\sigma_{\mathsf j}^2\le\mathsf v$ when $\mathsf c_{\mathsf j}=0$, and the budget binds because $\mathcal Y_{\mathsf j}'>0$. Solving gives $e^{2\mathsf c_{\mathsf j}}=\sigma_{\mathsf j}^2(1-\mathsf v)/(\mathsf v(1-\sigma_{\mathsf j}^2))$ for $\sigma_{\mathsf j}^2>\mathsf v$, then $\mathsf g_{\mathsf j}=1-e^{-2\mathsf c_{\mathsf j}}=(\sigma_{\mathsf j}^2-\mathsf v)/(\sigma_{\mathsf j}^2(1-\mathsf v))$ and $1-\mathsf g_{\mathsf j}\sigma_{\mathsf j}^2=(1-\sigma_{\mathsf j}^2)/(1-\mathsf v)$, which is the stated value. The capacity used is a continuous, strictly decreasing function of $\mathsf v\in(0,\sigma_1^2)$ that tends to infinity as $\mathsf v\to0$ and to zero as $\mathsf v\to\sigma_1^2$, so $\mathsf v$ is unique and decreases as $I_{\max}$ grows; variate $\mathsf m$ is active exactly when $\mathsf v<\sigma_{\mathsf m}^2$, that is, when the capacity used at $\mathsf v=\sigma_{\mathsf m}^2$, which is $\mathcal I_{\mathsf m}$, is below $I_{\max}$.

(c) The bound and its equality case are \cref{lem:env-eigen}(e). Under the budget, maximize $\sum_{\mathsf j}\lambda_{\mathsf j}(\mathsf J)(1-e^{-2\mathsf c_{\mathsf j}})$, again separable and concave: for active channels $2\lambda_{\mathsf j}(\mathsf J)e^{-2\mathsf c_{\mathsf j}}=2\mathsf v'$, so $\mathsf c_{\mathsf j}=\frac12\ln(\lambda_{\mathsf j}(\mathsf J)/\mathsf v')$, $\mathsf g_{\mathsf j}=1-\mathsf v'/\lambda_{\mathsf j}(\mathsf J)$, and the value is $\sum_{\mathsf j}(\lambda_{\mathsf j}(\mathsf J)-\mathsf v')^+$. The weighted error is \cref{lem:env-information}(iii).

(d) $\mathsf J_{\Sigma_{\mathsf Y}^{-1}}=\Sigma_{\mathsf X}^{-1/2}\Sigma_{\mathsf X\mathsf Y}\Sigma_{\mathsf Y}^{-1}\Sigma_{\mathsf Y\mathsf X}\Sigma_{\mathsf X}^{-1/2}=\mathsf C\mathsf C^\top$. Let $u_{\mathrm I}(\mathsf v)$ and $u_{\mathrm E}(\mathsf v)$ be the capacities that the information and the error allocations use at level $\mathsf v$, namely the sums over $\sigma_{\mathsf j}^2>\mathsf v$ of $\frac12\ln\frac{\sigma_{\mathsf j}^2(1-\mathsf v)}{\mathsf v(1-\sigma_{\mathsf j}^2)}$ and of $\frac12\ln\frac{\sigma_{\mathsf j}^2}{\mathsf v}$. Since $(1-\mathsf v)/(1-\sigma_{\mathsf j}^2)\ge1$ for $\sigma_{\mathsf j}^2\ge\mathsf v$, $u_{\mathrm I}\ge u_{\mathrm E}$ pointwise. Both decrease, so at a common budget $u_{\mathrm E}(\mathsf v_{\mathrm I})\le u_{\mathrm I}(\mathsf v_{\mathrm I})=u_{\mathrm E}(\mathsf v_{\mathrm E})$ forces $\mathsf v_{\mathrm I}\ge\mathsf v_{\mathrm E}$, and the information criterion's active set is contained in the error criterion's. For another weight, \cref{prop:env-horizon}(c) gives an instance in which $\mathsf J_I$ and $\mathsf C\mathsf C^\top$ are diagonal in one basis and their leading eigenvectors differ. Two symmetric matrices with a common orthonormal eigenbasis commute, so when $\mathsf J_\Upsilon\mathsf C\mathsf C^\top\neq\mathsf C\mathsf C^\top\mathsf J_\Upsilon$ no orthonormal eigenbasis of $\mathsf J_\Upsilon$ consists of eigenvectors of $\mathsf C\mathsf C^\top$, which are the left singular vectors of $\mathsf C$.

(e) For a noiseless sensor, $\mathsf G_s=\mathsf U\mathsf U^\top$ with $\mathsf U$ an orthonormal basis of the row space of $\tilde H$, and by Sylvester's identity $I(s;\mathsf Y)=-\frac12\ln\det(I_\kappa-\mathsf U^\top\mathsf C\mathsf C^\top\mathsf U)$. By \cref{lem:env-eigen}(a), $\lambda^\downarrow_{\mathsf j}(\mathsf U^\top\mathsf C\mathsf C^\top\mathsf U)\le\sigma_{\mathsf j}^2$, which gives the bound. Equality forces each factor $1-\lambda^\downarrow_{\mathsf j}(\mathsf U^\top\mathsf C\mathsf C^\top\mathsf U)\ge1-\sigma_{\mathsf j}^2>0$ to be an equality, so $\tr(\mathsf U^\top\mathsf C\mathsf C^\top\mathsf U)=\sum_{\mathsf j\le\kappa}\sigma_{\mathsf j}^2$, and \cref{lem:env-eigen}(b) under the gap identifies the span. The error statement is \cref{lem:env-information}(iii) with $\tr(\mathsf U^\top\mathsf J_\Upsilon\mathsf U)\le\sum_{\mathsf j\le\kappa}\lambda_{\mathsf j}(\mathsf J_\Upsilon)$.
\end{proof}

\begin{proof}[Proof of \cref{lem:env-markov}]
Let $Z=\E[\eta_t\mid\mathcal S]$, a linear function of $\mathcal S$. By the tower property $\E[\eta_t\mid Z]=Z$, so $\mathsf G_Z=\Cov(P_0^{-1/2}Z)$ satisfies $\Cov(P_0^{-1/2}\eta_t\mid Z)=I-\mathsf G_Z$, $0\preceq\mathsf G_Z\preceq I$ and $\rank\mathsf G_Z\le\kappa$. By the conditional independence of $\mathcal S$ and the Markov property, $\E[\tilde{\mathsf Y}\mid\mathcal S]=\E\bigl[\E[\tilde{\mathsf Y}\mid(\eta_{t'})_{t'\le t}]\bigm|\mathcal S\bigr]=\E\bigl[\E[\tilde{\mathsf Y}\mid\eta_t]\bigm|\mathcal S\bigr]=\mathsf C^\top P_0^{-1/2}Z$, so $\Cov(\tilde{\mathsf Y}\mid\mathcal S)=I-\mathsf C^\top\mathsf G_Z\mathsf C$ and $I(\mathcal S;\mathsf Y)=-\frac12\ln\det(I-\mathsf C^\top\mathsf G_Z\mathsf C)$. Data processing gives $I(Z;\eta_t)\le I(\mathcal S;\eta_t)\le I(\mathcal S;(\eta_{t'})_{t'\le t})\le I_{\max}$, and $I(Z;\eta_t)=-\frac12\ln\det(I-\mathsf G_Z)$ is finite, so $\mathsf G_Z\prec I$. By \cref{lem:env-information}(iv) a sensor of $\eta_t$ with captured covariance $\mathsf G_Z$ exists; it has the same information about $\mathsf Y$ as $\mathcal S$, uses at most $I_{\max}$, and \cref{thm:env-canonical}(b) bounds it.
\end{proof}

\begin{proof}[Proof of \cref{prop:env-horizon}]
(a) Detailed balance gives $\mathsf F^\top=P_0^{1/2}F^\top P_0^{-1/2}=P_0^{-1/2}(P_0F^\top)P_0^{-1/2}=P_0^{-1/2}FP_0^{1/2}=\mathsf F$, and $\Cov(\eta_t,\eta_{t+h})=P_0(F^h)^\top$ gives $\mathsf C(h)=(\mathsf F^h)^\top=\mathsf F^h$. With $\mathsf F=\sum_{\mathsf j}\lambda_{\mathsf j}(F)\mathsf e_{\mathsf j}\mathsf e_{\mathsf j}^\top$, the singular values of $\mathsf F^h$ are $\lvert\lambda_{\mathsf j}(F)\rvert^h$ with singular vectors $\mathsf e_{\mathsf j}$, and $\mathsf u_{\mathsf j}=P_0^{-1/2}\mathsf e_{\mathsf j}$ satisfies $\mathsf u_{\mathsf j}^\top F=\mathsf e_{\mathsf j}^\top\mathsf FP_0^{-1/2}=\lambda_{\mathsf j}(F)\mathsf u_{\mathsf j}^\top$ and $\mathsf u_{\mathsf j}^\top P_0\mathsf u_{\mathsf j}=1$. The noiseless statement is \cref{thm:env-canonical}(e), and the activation condition is \cref{thm:env-canonical}(b). With $\sigma_{\mathsf j}(h)=e^{-h/\tau_{\mathsf j}}$, the $\mathsf j$-th term of $\mathcal I_{\mathsf m}(h)$ is $h(\tau_{\mathsf m}^{-1}-\tau_{\mathsf j}^{-1})+\frac12\ln\frac{1-e^{-2h/\tau_{\mathsf m}}}{1-e^{-2h/\tau_{\mathsf j}}}$, whose derivative in $h$ is $\mathcal Z(\tau_{\mathsf m}^{-1})-\mathcal Z(\tau_{\mathsf j}^{-1})$ with $\mathcal Z(y)=y/(1-e^{-2yh})=(2h)^{-1}x/(1-e^{-x})$, $x=2yh$. The function $x/(1-e^{-x})$ increases on $x>0$ (its derivative has numerator $1-e^{-x}(1+x)>0$), and $\tau_{\mathsf m}^{-1}\ge\tau_{\mathsf j}^{-1}$ for $\mathsf j<\mathsf m$, so each term is non-decreasing in $h$. The logarithm is non-negative because $\tau_{\mathsf m}\le\tau_{\mathsf j}$, which gives the linear lower bound; the last claim is the case $\mathsf m=2$.

(b) Write $\mathsf F^h=\mathsf M_h+\mathsf N_h$ with $\mathsf M_h=P_0^{-1/2}F^h\mathbb P_{\mathrm{slow}}P_0^{1/2}$ and $\mathsf N_h=P_0^{-1/2}F^h(I-\mathbb P_{\mathrm{slow}})P_0^{1/2}$, so that $\Vert\mathsf N_h\Vert\le\mathfrak a_h$ and $\sigma_\kappa(\mathsf M_h)\ge\mathfrak b_h$. A top-$\kappa$ canonical subspace of $\mathsf X$ is spanned by top-$\kappa$ left singular vectors of $\mathsf C(h)=(\mathsf F^h)^\top$, that is, by top-$\kappa$ right singular vectors of $\mathsf F^h$, the columns of some $\mathsf Z_\kappa$ with $(\mathsf F^h)^\top\mathsf W_\kappa=\mathsf Z_\kappa\mathsf T_\kappa$, where $\mathsf W_\kappa$ holds the left singular vectors and $\mathsf T_\kappa$ the top $\kappa$ singular values. Since $F^h\mathbb P_{\mathrm{slow}}=\mathbb P_{\mathrm{slow}}F^h$ and the slow eigenvalues are nonzero, the row space of $F^h\mathbb P_{\mathrm{slow}}$ is $\mathcal U_\kappa$ and that of $\mathsf M_h$ is $P_0^{1/2}\mathcal U_\kappa$. Let $\mathbb P_\perp$ project orthogonally onto the complement of $P_0^{1/2}\mathcal U_\kappa$, so $\mathbb P_\perp\mathsf M_h^\top=0$. Then $\mathbb P_\perp\mathsf Z_\kappa=\mathbb P_\perp\mathsf N_h^\top\mathsf W_\kappa\mathsf T_\kappa^{-1}$, and Weyl's inequality $\sigma_\kappa(\mathsf F^h)\ge\sigma_\kappa(\mathsf M_h)-\Vert\mathsf N_h\Vert$ gives $\sin\vartheta_h=\Vert\mathbb P_\perp\mathsf Z_\kappa\Vert\le\mathfrak a_h/(\mathfrak b_h-\mathfrak a_h)$. (Since $\sigma_{\kappa+1}(\mathsf F^h)\le\Vert\mathsf N_h\Vert$, the subspace is unique once $\mathfrak a_h<\mathfrak b_h/2$.) With $F=\mathsf Q\diag(\lambda(F))\mathsf Q^{-1}$, the slow and fast columns $\mathsf Q_{\mathrm s},\mathsf Q_{\mathrm f}$ of the eigenvector matrix $\mathsf Q$ and rows $\mathsf U_{\mathrm s},\mathsf U_{\mathrm f}$ of $\mathsf Q^{-1}$ give $F^h(I-\mathbb P_{\mathrm{slow}})=\mathsf Q_{\mathrm f}\diag(\lambda_{\mathrm{fast}}(F))^h\mathsf U_{\mathrm f}$ and $F^h\mathbb P_{\mathrm{slow}}=\mathsf Q_{\mathrm s}\diag(\lambda_{\mathrm{slow}}(F))^h\mathsf U_{\mathrm s}$, hence
\begin{equation}\label{eq:env-ab-bounds}
\mathfrak a_h\le\mathsf c_P\Vert\mathsf Q\Vert\Vert\mathsf Q^{-1}\Vert\lvert\lambda_{\kappa+1}(F)\rvert^h,\qquad\mathfrak b_h\ge\sigma_\kappa(\mathsf Q_{\mathrm s})\sigma_\kappa(\mathsf U_{\mathrm s})\lvert\lambda_\kappa(F)\rvert^h/\mathsf c_P,
\end{equation}
which gives the rate. A canonical direction $z\in P_0^{1/2}\mathcal U_\kappa$ is the functional $z^\top P_0^{-1/2}\eta_t$ with $P_0^{-1/2}z\in\mathcal U_\kappa$, a combination of slow mode amplitudes. For the error criterion the optimal readings are the top right singular vectors of $\Upsilon^{1/2}F^hP_0^{1/2}$ (\cref{thm:env-canonical}(e)), and the same splitting applies.

(c) Here $P_0=Q(I-F^2)^{-1}$ shares the eigenbasis of $F$, $\mathsf J_I=P_0^{1/2}(F^h)^\top F^hP_0^{1/2}$ is diagonal with entries $\lambda_{\mathsf j}(F)^{2h}\mathsf p_{\mathsf j}$, and \cref{thm:env-canonical}(e) applies, so the rank-one error optimum reads a mode of largest entry. The ratio of the entries of the fast and the slow mode, $(\mathsf p_{\mathrm f}/\mathsf p_{\mathrm s})(\lvert\lambda_{\mathrm f}(F)\rvert/\lvert\lambda_{\mathrm s}(F)\rvert)^{2h}$, exceeds one at $h=0$, decreases strictly in $h$ and equals one at $h^\ast$, which gives the ranking of the two modes on either side of $h^\ast$. The canonical correlations are $\lvert\lambda_{\mathsf j}(F)\rvert^h$, larger for the slow mode than for the fast mode at every $h\ge1$. With only these two modes the larger entry and the larger correlation are the maxima, which gives the modes read. With a third mode neither maximum need lie in the pair: for $\lambda(F)=(0.97,0.9,0.6)$ and $\mathsf p=(0.3384,5,1.5625)$ the entries at $h=1$ are $(0.318,4.05,0.5625)$, and the error optimum reads the second mode, and a third mode with $\lvert\lambda_{\mathsf j}(F)\rvert>\lvert\lambda_{\mathrm s}(F)\rvert$ has the largest canonical correlation at every $h\ge1$.
\end{proof}

\begin{proof}[Proof of \cref{prop:env-diffusive}]
$F$ is symmetric and commutes with $Q=Q_0I$, so $P_0=Q_0(I-F^2)^{-1}$ commutes with $F$ and $FP_0=P_0F^\top$. The eigenvalues $1-\varepsilon_{\mathrm w}-\mathsf d_{\mathrm w}\lambda(L_{\mathrm w})$ of $F$ lie in $(0,1)$ and decrease with the Laplacian eigenvalue, so the order by modulus is the order by frequency, with the constant vector, of Laplacian eigenvalue zero, first. \Cref{prop:env-horizon}(a) gives the rest; repeated Laplacian eigenvalues do no harm, since modes of equal modulus share one threshold and the terms of $\mathcal I_{\mathsf m}(h)$ that pair them vanish, and $\tau_2<\tau_1$ because the graph is connected.
\end{proof}

\begin{proof}[Proof of \cref{thm:env-control}]
(a) Write $\hat\eta=\E[\eta_{t+h}\mid s]$ and $\Sigma_s=\Cov(\eta_{t+h}\mid s)$. Given $s$, $o-o^\ast=(H^{\mathrm o}\hat\eta+B^{\mathrm o}a-o^\ast)+H^{\mathrm o}(\eta_{t+h}-\hat\eta)+v^{\mathrm o}$, and the last two terms have conditional mean zero, so
\[
2\E[-\ln\tilde p(o)\mid s]-\ln\det(2\pi\Psi_{\mathrm o})=\bigl\Vert H^{\mathrm o}\hat\eta+B^{\mathrm o}a-o^\ast\bigr\Vert^2_{\Psi_{\mathrm o}^{-1}}+\tr\bigl(\Psi_{\mathrm o}^{-1}(H^{\mathrm o}\Sigma_s(H^{\mathrm o})^\top+R^{\mathrm o})\bigr).
\]
Weighted least squares minimizes the first term by $B^{\mathrm o}a^\ast=-\mathbb P_{\mathrm o}(H^{\mathrm o}\hat\eta-o^\ast)$, which is the stated $a^\ast$, with residual $\Vert(I-\mathbb P_{\mathrm o})(H^{\mathrm o}\hat\eta-o^\ast)\Vert^2_{\Psi_{\mathrm o}^{-1}}$; since the act may depend on $s$ arbitrarily, the minimal expectation is the expectation of this conditional minimum. The matrix $\Psi_{\mathrm o}^{-1}\mathbb P_{\mathrm o}$ is symmetric and $\mathbb P_{\mathrm o}$ is idempotent, so $\mathbb P_{\mathrm o}$ is an orthogonal projector in the $\Psi_{\mathrm o}^{-1}$ inner product, $\Vert x\Vert^2=\Vert\mathbb P_{\mathrm o}x\Vert^2+\Vert(I-\mathbb P_{\mathrm o})x\Vert^2$ in that norm, and $\mathbb P_{\mathrm o}^\top\Psi_{\mathrm o}^{-1}\mathbb P_{\mathrm o}=\Psi_{\mathrm o}^{-1}\mathbb P_{\mathrm o}\succeq0$. Since $\E\hat\eta=0$ and $\Cov(\hat\eta)=P_0-\Sigma_s$, the expected residual is a constant plus $\tr\bigl(\Psi_{\mathrm o}^{-1}(I-\mathbb P_{\mathrm o})H^{\mathrm o}(P_0-\Sigma_s)(H^{\mathrm o})^\top(I-\mathbb P_{\mathrm o})^\top\bigr)$, and adding $\tr(\Psi_{\mathrm o}^{-1}H^{\mathrm o}\Sigma_s(H^{\mathrm o})^\top)$ leaves, by the Pythagorean identity, the sensor-dependent part $\tr(\Psi_{\mathrm o}^{-1}\mathbb P_{\mathrm o}H^{\mathrm o}\Sigma_s(H^{\mathrm o})^\top\mathbb P_{\mathrm o}^\top)=\tr(\Upsilon_{\mathrm{ctl}}\Sigma_s)$. The rank bound is $\rank\mathbb P_{\mathrm o}=\rank B^{\mathrm o}$. With no sensor $\Sigma_s=P_0$, and $P_0-\Sigma_s=F^hP_0^{1/2}\mathsf G_sP_0^{1/2}(F^h)^\top$ by the identity in the proof of \cref{lem:env-information}(iii) with $\Sigma_{\mathsf X\mathsf Y}=P_0(F^h)^\top$, which gives the pragmatic value $\frac12\tr(\mathsf G_s\mathsf J_{\mathrm{ctl}})$. Finally $P_0-\Sigma_s=\Cov(\eta_{t+h},s)\Cov(s)^{-1}\Cov(s,\eta_{t+h})$, so $\tr(\Upsilon_{\mathrm{ctl}}(P_0-\Sigma_s))=\tr(\Psi_{\mathrm o}^{-1}KK^\top)$ with $K=\Cov(\mathbb P_{\mathrm o}H^{\mathrm o}\eta_{t+h},s)\Cov(s)^{-1/2}$, which vanishes if and only if $K=0$.

(b) \Cref{thm:env-canonical}(c) with $\mathsf J=\mathsf J_{\mathrm{ctl}}$, whose rank is at most $\rank\Upsilon_{\mathrm{ctl}}\le\rank B^{\mathrm o}$; reverse water-filling gives capacity only to positive eigenvalues.

(c) By (a), $G^{\mathrm{att}}_h=\mathfrak c_0+\frac12\tr(\Upsilon_{\mathrm{ctl}}P_0)-I(s;\mathsf S_{\mathrm m}\eta_{t+h})-\frac12\tr(\mathsf G_s\mathsf J_{\mathrm{ctl}})$. By \cref{thm:env-canonical}(a) and \cref{lem:env-eigen}(e), $-I(s;\mathsf S_{\mathrm m}\eta_{t+h})\ge\sum_{\mathsf j\le\kappa}\frac12\ln(1-\mathsf g_{\mathsf j}\sigma_{\mathsf j}^2)$ and $-\frac12\tr(\mathsf G_s\mathsf J_{\mathrm{ctl}})\ge-\frac12\sum_{\mathsf j\le\kappa}\mathsf g_{\mathsf j}\lambda_{\mathsf j}(\mathsf J_{\mathrm{ctl}})$, both with the sorted captured fractions; under the hypothesis the sensor whose $\mathsf j$-th channel reads $\mathsf e_{\mathsf j}$ attains both. Minimizing the right side over the captured fractions is separable in the capacities $\mathsf c_{\mathsf j}=-\frac12\ln(1-\mathsf g_{\mathsf j})$, with per-channel marginal value
\[
\frac{\sigma_{\mathsf j}^2e^{-2\mathsf c_{\mathsf j}}}{1-\sigma_{\mathsf j}^2+\sigma_{\mathsf j}^2e^{-2\mathsf c_{\mathsf j}}}+\lambda_{\mathsf j}(\mathsf J_{\mathrm{ctl}})e^{-2\mathsf c_{\mathsf j}},
\]
which decreases in $\mathsf c_{\mathsf j}$ and does not decrease in $(\sigma_{\mathsf j}^2,\lambda_{\mathsf j}(\mathsf J_{\mathrm{ctl}}))$. The objective is concave in the capacities, so the Karush--Kuhn--Tucker conditions characterize the optimum, and if channel $\mathsf j+1$ is active with marginal value equal to the multiplier, channel $\mathsf j$, whose marginal value function lies above, receives at least as much capacity. The active set is therefore a leading segment with non-increasing capacities; the same argument applies to each of the three objectives.
\end{proof}

\begin{proof}[Proof of \cref{prop:env-selection}]
The value splits as stated by the chain rule $I(s^{\mathrm{old}},s^{(c)};\mathsf S_{\mathrm m}\eta_{t+h})=I(s^{\mathrm{old}};\mathsf S_{\mathrm m}\eta_{t+h})+\mathsf V_{\mathrm{ep}}$ and by the identity $\min_{a(\cdot)}\E[-\ln\tilde p(o)]=\mathfrak c_0+\frac12\tr(\Upsilon_{\mathrm{ctl}}\Cov(\eta_{t+h}\mid s))$ of the proof of \cref{thm:env-control}(a), whose derivation uses only that the readings $s$ are jointly Gaussian with $\eta_{t+h}$ and independent of $v^{\mathrm o}$. (i) For jointly Gaussian variables a vanishing conditional covariance is conditional independence, so $s^{(c)}$ is independent of $\eta_{t+h}$ given $s^{\mathrm{old}}$; the conditional information vanishes and the conditional covariance of $\eta_{t+h}$ does not change. Both increments are continuous functions of the joint covariance, which is continuous in $H^{(c)}$, because the nondegenerate noise keeps $\Cov(s^{(c)}\mid s^{\mathrm{old}})$ invertible. (ii) As in the proof of \cref{thm:env-control}(a), conditionally on $s^{\mathrm{old}}$, $2\mathsf V_{\mathrm{pr}}=\tr(\Psi_{\mathrm o}^{-1}KK^\top)$ with $K=\Cov(\mathbb P_{\mathrm o}H^{\mathrm o}\eta_{t+h},s^{(c)}\mid s^{\mathrm{old}})\Cov(s^{(c)}\mid s^{\mathrm{old}})^{-1/2}$, zero if and only if $K=0$; the epistemic part is Gaussian conditional independence. (iii) By independence $I(s^{(c)};s^{\mathrm{old}})=0$, so $\mathsf V_{\mathrm{ep}}=I(s^{(c)};\mathsf S_{\mathrm m}\eta_{t+h},s^{\mathrm{old}})\le I(\eta^{(c)}_t;\eta_{t+h},s^{\mathrm{old}})=I(\eta^{(c)}_t;\eta^{(c)}_{t+h})$, using data processing and the fact that $(\eta_{t+h},s^{\mathrm{old}})$ consists of $\eta^{(c)}_{t+h}$ and variables independent of content $c$; the last quantity is the noiseless value $-\frac12\sum_{\mathsf j}\ln(1-\sigma^{(c)}_{\mathsf j}(h)^2)$ of the block. For the pragmatic part, independence gives $\Cov(\eta_{t+h},s^{(c)}\mid s^{\mathrm{old}})=\Cov(\eta_{t+h},s^{(c)})$ and $\Cov(s^{(c)}\mid s^{\mathrm{old}})=\Cov(s^{(c)})$, so the reduction of the covariance is $\Cov(\E[\eta_{t+h}\mid s^{(c)}])\preceq\Cov(\E[\eta_{t+h}\mid\eta^{(c)}_t])$, the embedding of $F^{(c)h}P^{(c)}_0(F^{(c)h})^\top$, and $\tr(\Upsilon_{\mathrm{ctl}}X)\le\Vert\Upsilon_{\mathrm{ctl}}\Vert\tr X$ for $X\succeq0$. The bound on $\sigma^{(c)}_{\mathsf j}(h)$ is the norm of $(P^{(c)}_0)^{1/2}(F^{(c)h})^\top(P^{(c)}_0)^{-1/2}$, which vanishes geometrically because $\rho(F^{(c)})<1$: by the independence of the contents and $P_0\succ0$, $F^{(c)h}$ is the diagonal block of $F^h$ for content $c$, so $\rho(F^{(c)})\le\rho(F)<1$. (iv) is \cref{lem:env-constant}.
\end{proof}

\begin{proof}[Proof of \cref{lem:env-constant}]
The predicted variance obeys $P^-_{t+1}=f(P^-_t)$ with $f(x)=F_{\mathrm g}^2xR^s/(x+R^s)+Q_{\mathrm g}$, increasing and concave on $[0,\infty)$ with $f(0)=Q_{\mathrm g}>0$ and $f(x)-x\to-\infty$. Hence $f(x)-x$ has exactly one positive root, $f(x)>x$ below it and $f(x)<x$ above it, and the iteration converges monotonically from any prior. The fixed-point equation $x(x+R^s)=F_{\mathrm g}^2xR^s+Q_{\mathrm g}(x+R^s)$ is the stated quadratic, whose positive root $\frac12\bigl(-(\mathsf b-Q_{\mathrm g})+((\mathsf b-Q_{\mathrm g})^2+4Q_{\mathrm g}R^s)^{1/2}\bigr)$, $\mathsf b=R^s(1-F_{\mathrm g}^2)$, does not increase in $\mathsf b$; at $\mathsf b=0$ it is $\frac12(Q_{\mathrm g}+(Q_{\mathrm g}^2+4Q_{\mathrm g}R^s)^{1/2})\le Q_{\mathrm g}+(Q_{\mathrm g}R^s)^{1/2}$. The innovation $s_t-\E[s_t\mid s_{<t}]$ has variance $P^-_t+R^s$, and $R^s$ given $\eta^{\mathrm g}_t$, which gives the information. Given the past readings, $s_t\to\eta^{\mathrm g}_t\to\eta^{\mathrm g}_{t+h}$ is a Markov chain, so data processing bounds the information about the future. For a constant, the predicted variance after $n_{\mathrm r}$ readings is $(1/P_{\mathrm{pr}}+n_{\mathrm r}/R^s)^{-1}$, and the total information is one half the logarithm of the ratio of $P_{\mathrm{pr}}$ to it.
\end{proof}

\begin{proof}[Proof of \cref{prop:env-concentration}]
Let $\mathsf k_u(z)=(4\pi D_{\mathrm m}u)^{-3/2}e^{-\lvert z\rvert^2/(4D_{\mathrm m}u)}$ be the heat kernel, and $\mathfrak K(u)=\int_{\mathbb B}\int_{\mathbb B}\mathsf k_u(x-y)\mathrm dx\mathrm dy$ for $u>0$, $\mathfrak K(0)=\lvert\mathbb B\rvert$. Restrict first to the particles that start in a ball $B_{R_0}$ of radius $R_0$: their number $\mathsf N_0$ is Poisson with mean $\bar c_{\mathrm L}\lvert B_{R_0}\rvert$, each starts uniformly in $B_{R_0}$, and their paths are independent. For bounded functionals $\mathsf a_{\mathsf m},\mathsf b_{\mathsf m}$ of the path of particle $\mathsf m$, $\E[\sum_{\mathsf m}\mathsf a_{\mathsf m}\sum_{\mathsf m}\mathsf b_{\mathsf m}]=\E[\mathsf N_0]\E[\mathsf a_1\mathsf b_1]+\E[\mathsf N_0(\mathsf N_0-1)]\E\mathsf a_1\E\mathsf b_1$ with $\E[\mathsf N_0(\mathsf N_0-1)]=\E[\mathsf N_0]^2$, so the covariance of the sums is $\E[\mathsf N_0]\E[\mathsf a_1\mathsf b_1]$. With $\mathsf a$ and $\mathsf b$ the indicators that the particle lies in a set $A$ at time $t$ and in $A'$ at time $t'\ge t$, and $R_0\to\infty$,
\[
\Cov\bigl(\mathsf N_A(t),\mathsf N_{A'}(t')\bigr)=\bar c_{\mathrm L}\int\Bigl[\int_A\mathsf k_t(y-x)\int_{A'}\mathsf k_{t'-t}(y'-y)\mathrm dy'\mathrm dy\Bigr]\mathrm dx=\bar c_{\mathrm L}\int_A\int_{A'}\mathsf k_{t'-t}(y'-y)\mathrm dy'\mathrm dy.
\]
Integrating over $[0,T_{\mathrm{int}}]^2$ gives $\Var\hat c=2\bar c_{\mathrm L}(\lvert\mathbb B\rvert T_{\mathrm{int}})^{-2}\int_0^{T_{\mathrm{int}}}(T_{\mathrm{int}}-u)\mathfrak K(u)\mathrm du$, and $\E\hat c=\bar c_{\mathrm L}$ is stationarity. Leading term: $\int_0^\infty\mathsf k_u(z)\mathrm du=1/(4\pi D_{\mathrm m}\lvert z\rvert)$, so $\int_0^\infty\mathfrak K=(4\pi D_{\mathrm m})^{-1}\int_{\mathbb B}\int_{\mathbb B}\lvert x-y\rvert^{-1}$. For a function $\mathsf f$ of the distance, $\int_{\mathbb B}\int_{\mathbb B}\mathsf f(\lvert x-y\rvert)=\int_0^{2r_{\mathrm{cell}}}4\pi z^2V_\cap(z)\mathsf f(z)\mathrm dz$, with the lens volume $V_\cap(z)=\frac\pi{12}(4r_{\mathrm{cell}}+z)(2r_{\mathrm{cell}}-z)^2$; with $z=2r_{\mathrm{cell}}\vartheta$,
\[
\int_0^{2r_{\mathrm{cell}}}zV_\cap(z)\mathrm dz=\frac{32\pi r_{\mathrm{cell}}^5}{12}\int_0^1\vartheta(2+\vartheta)(1-\vartheta)^2\mathrm d\vartheta=\frac{8\pi r_{\mathrm{cell}}^5}{15},
\]
so $\int_{\mathbb B}\int_{\mathbb B}\lvert x-y\rvert^{-1}=32\pi^2r_{\mathrm{cell}}^5/15=6\lvert\mathbb B\rvert^2/(5r_{\mathrm{cell}})$ and $2\bar c_{\mathrm L}(\lvert\mathbb B\rvert^2T_{\mathrm{int}})^{-1}\int_0^\infty\mathfrak K=3\bar c_{\mathrm L}/(5\pi D_{\mathrm m}r_{\mathrm{cell}}T_{\mathrm{int}})$. Corrections: since $\lvert x-y\rvert\le2r_{\mathrm{cell}}$ on $\mathbb B\times\mathbb B$, $\exp(-\lvert x-y\rvert^2/(4D_{\mathrm m}u))\ge1-r_{\mathrm{cell}}^2/(D_{\mathrm m}u)$, which gives
\[
\frac{\lvert\mathbb B\rvert^2}{(4\pi D_{\mathrm m}u)^{3/2}}\Bigl(1-\frac{r_{\mathrm{cell}}^2}{D_{\mathrm m}u}\Bigr)\le\mathfrak K(u)\le\frac{\lvert\mathbb B\rvert^2}{(4\pi D_{\mathrm m}u)^{3/2}}.
\]
Write
\[
\Var\hat c=\frac{2\bar c_{\mathrm L}}{\lvert\mathbb B\rvert^2T_{\mathrm{int}}}\Bigl[\int_0^\infty\mathfrak K-\int_{T_{\mathrm{int}}}^\infty\mathfrak K\Bigr]-\frac{2\bar c_{\mathrm L}}{\lvert\mathbb B\rvert^2T_{\mathrm{int}}^2}\int_0^{T_{\mathrm{int}}}u\mathfrak K(u)\mathrm du.
\]
The bounds give, splitting the second integral at $u=r_{\mathrm{cell}}^2/D_{\mathrm m}$ and using the lower bound only above that point,
\[
\int_{T_{\mathrm{int}}}^\infty\mathfrak K=\frac{2\lvert\mathbb B\rvert^2}{(4\pi D_{\mathrm m})^{3/2}T_{\mathrm{int}}^{1/2}}-\mathrm{rem}_2,\qquad\int_0^{T_{\mathrm{int}}}u\mathfrak K(u)\mathrm du=\frac{2\lvert\mathbb B\rvert^2T_{\mathrm{int}}^{1/2}}{(4\pi D_{\mathrm m})^{3/2}}-\mathrm{rem}_3,
\]
with
\[
0\le\mathrm{rem}_2\le\frac{2\lvert\mathbb B\rvert^2r_{\mathrm{cell}}^2}{3(4\pi D_{\mathrm m})^{3/2}D_{\mathrm m}T_{\mathrm{int}}^{3/2}},\qquad0\le\mathrm{rem}_3\le\frac{4\lvert\mathbb B\rvert^2r_{\mathrm{cell}}}{(4\pi D_{\mathrm m})^{3/2}D_{\mathrm m}^{1/2}}.
\]
The two terms of order $T_{\mathrm{int}}^{-3/2}$ sum to $-8\bar c_{\mathrm L}(4\pi D_{\mathrm m}T_{\mathrm{int}})^{-3/2}$, which relative to the leading term is $-\frac53(r_{\mathrm{cell}}^2/(\pi D_{\mathrm m}T_{\mathrm{int}}))^{1/2}$, and $\mathrm{rem}_2$ and $\mathrm{rem}_3$ enter with positive sign and give the two parts of the bound on $\mathfrak q_T$. In an ideal solution the chemical potential is a constant plus $k_{\mathrm B}T_{\mathrm K}$ times the logarithm of the concentration, so its small changes in units of $k_{\mathrm B}T_{\mathrm K}$ are relative changes of concentration.
\end{proof}

\begin{proof}[Proof of \cref{prop:env-aggregate-closure}]
(a) By \cref{def:env-aggregates}, $\bar x$ is exactly closed if and only if $\ker\one^\top$ is invariant under $W$, that is, $\one^\top Wx=0$ for every $x\perp\one$, that is, $W^\top\one\in\operatorname{span}(\one)$; then $N^{-1}\one^\top Wx=\bar w\bar x$. (b) Averaging the agents' equation over type $\mathfrak t$ gives the type equation, because the coupling depends on the other agents only through type averages and $A_{\mathfrak t}$ is common to the type; the averaged idiosyncratic noise is independent of the past, with covariance $Q_{\mathrm{id}}/N_{\mathfrak t}$, and the world is autonomous. The deviation from the infinite-population limit obeys $\bar x^{\mathrm{fl}}_{t+1}=\mathbb A_{\mathrm{ty}}\bar x^{\mathrm{fl}}_t+\bar w^{\mathrm{id}}_t$, driven by noise independent of the world and of the common shocks, which gives $\Sigma^{\mathrm{fl}}$ and its bound. (c) For agents $i,i'$ of type $\mathfrak t$ with own dynamics $A_i\ne A_{i'}$, the state with $x_i=e$, $x_{i'}=-e$ and all other coordinates zero lies in the kernel of the map to type averages, and its image has type-$\mathfrak t$ average $(A_i-A_{i'})e/N_{\mathfrak t}$, nonzero for a suitable $e$; the kernel is not invariant.
\end{proof}

\begin{proof}[Proof of \cref{prop:env-aggregate-memory}]
(a) If $\mathsf R(u)=\mathsf R(0)\alpha_\ast^{\lvert u\rvert}$, then $\Cov(x_{t+1}-\alpha_\ast x_t,x_{t-l})=\mathsf R(l+1)-\alpha_\ast\mathsf R(l)=0$ for $l\ge0$, so the Gaussian residual is independent of the past and the process is an autoregression of order one. If the process is Markov, $\E[x_{t+u}\mid x_{\le t}]=(\mathsf R(u)/\mathsf R(0))x_t$, and the tower property applied to $\E[x_{t+u+1}\mid x_{\le t}]$ gives $\mathsf R(u+1)/\mathsf R(0)=(\mathsf R(1)/\mathsf R(0))(\mathsf R(u)/\mathsf R(0))$, so the autocovariance is geometric. A geometric autocorrelation is $\Corr(\bar x_t,\bar x_{t+1})^u$.

(b) The stationary solution is $x_{i,t}=\sum_{l\ge0}\alpha_i^l(w^{\mathrm{com}}_{t-1-l}+w^{\mathrm{id}}_{i,t-1-l})$, so $\bar x_t=\sum_l\bar{\mathfrak m}_lw^{\mathrm{com}}_{t-1-l}+N^{-1}\sum_i\sum_l\alpha_i^lw^{\mathrm{id}}_{i,t-1-l}$ with $\bar{\mathfrak m}_l=N^{-1}\sum_i\alpha_i^l$, and
\[
\mathsf R(u)=q_{\mathrm{com}}\sum_{l\ge0}\bar{\mathfrak m}_l\bar{\mathfrak m}_{l+u}+\frac{q_{\mathrm{id}}}{N^2}\sum_i\frac{\alpha_i^u}{1-\alpha_i^2},\qquad\sum_{l\ge0}\bar{\mathfrak m}_l\bar{\mathfrak m}_{l+u}=\frac1{N^2}\sum_{i,i'}\frac{\alpha_{i'}^u}{1-\alpha_i\alpha_{i'}},
\]
which is the stated measure; its density with respect to $\mathbb Q$ is positive because $\lvert\alpha\alpha'\rvert<1$.

(c) By (a), $\bar x$ is Markov if and only if $\int\alpha^u\bar{\mathbb M}(\mathrm d\alpha)=\alpha_\ast^u$ for all $u\ge0$ and some $\alpha_\ast$, with $\bar{\mathbb M}=\mathbb M/\mathbb M(\R)$; the probability measure $\bar{\mathbb M}$ on $[-1,1]$ then has the moments of the point mass at $\alpha_\ast$, and moments determine such a measure, so $\bar{\mathbb M}$ is that point mass. Since $\mathbb M$ charges every $\alpha_i$ (as $q_{\mathrm{com}}+q_{\mathrm{id}}>0$), all $\alpha_i$ equal $\alpha_\ast$. Equal persistences give an autoregression of order one.

(d) With $\bar\alpha=\int\alpha\bar{\mathbb M}(\mathrm d\alpha)=\Corr(\bar x_t,\bar x_{t+1})$ and $\varphi(\alpha)=\alpha^u$, convex on $[0,1)$, Taylor's theorem gives
\[
0\le\int\varphi\mathrm d\bar{\mathbb M}-\varphi(\bar\alpha)=\int\bigl[\varphi(\alpha)-\varphi(\bar\alpha)-\varphi'(\bar\alpha)(\alpha-\bar\alpha)\bigr]\bar{\mathbb M}(\mathrm d\alpha)\le\tfrac12u(u-1)\alpha_+^{u-2}\Var_{\bar{\mathbb M}}(\alpha),
\]
for $u\ge2$ (the error vanishes for $u\le1$), and $\frac12u(u-1)x^{u-2}\le\sum_{\mathsf m\ge2}\frac12\mathsf m(\mathsf m-1)x^{\mathsf m-2}=(1-x)^{-3}$. A law on an interval of length $\alpha_+-\alpha_-$ has variance at most a quarter of its square.

(e) The type averages obey (b) of \cref{prop:env-aggregate-closure} with $\mathsf B=0$ and $A_{\mathfrak t}$ the common persistence, and $\bar x$ is a fixed linear combination of them; adjoining $\mathfrak T-1$ of them to $\bar x$ spans all, so $\bar x$ is closed of order $\mathfrak T-1$. No lower order suffices. By \cref{thm:met-abstract-physics}(f), applied without input as in \cref{def:env-aggregates}, every exactly closed extension of $\bar x$ has a row space that contains the row space of the observability matrix $\mathcal O$ of $(N^{-1}\one^\top,\diag(\alpha_i))$, whose rows are $N^{-1}(\alpha_1^l,\dots,\alpha_N^l)$ for $0\le l<N$. The columns of $\mathcal O$ that belong to agents of one type coincide, and its $\mathfrak T$ distinct columns are those of a Vandermonde matrix with distinct nodes, multiplied by $N^{-1}$, so $\rank\mathcal O=\mathfrak T$ and $d^\ast=\rank\mathcal O-1=\mathfrak T-1$.

(f) Put $\alpha=1-y/l$: $l^{\mathfrak p}\mathfrak m_l=\int_0^l(1-y/l)^ly^{\mathfrak p-1}\mathsf f_0(1-y/l)\mathrm dy\to\mathsf f_0(1)\Gamma(\mathfrak p)$ by dominated convergence, with dominating function $\Vert\mathsf f_0\Vert_\infty e^{-y}y^{\mathfrak p-1}$; the same bound gives $\mathfrak m_l\le\Vert\mathsf f_0\Vert_\infty\Gamma(\mathfrak p)l^{-\mathfrak p}$ for $l\ge1$. Since $\sum_l\mathfrak m_l^2<\infty$, the limiting aggregate $\sum_l\mathfrak m_lw^{\mathrm{com}}_{t-1-l}$ exists in mean square, with $\mathsf R(u)/q_{\mathrm{com}}=\sum_{l\ge0}\mathfrak m_l\mathfrak m_{l+u}$. The term $l=0$ is $\mathfrak m_u=O(u^{-\mathfrak p})=o(u^{1-2\mathfrak p})$ because $\mathfrak p<1$. For $l\ge1$, $u^{2\mathfrak p-1}\mathfrak m_l\mathfrak m_{l+u}=\int_{(l-1)/u}^{l/u}\mathsf h_u(x)\mathrm dx$ with the step function $\mathsf h_u(x)=u^{\mathfrak p}\mathfrak m_{\lceil ux\rceil}u^{\mathfrak p}\mathfrak m_{\lceil ux\rceil+u}$, which converges pointwise to $(\mathsf f_0(1)\Gamma(\mathfrak p))^2x^{-\mathfrak p}(1+x)^{-\mathfrak p}$ and, because $\lceil ux\rceil\ge ux$, is dominated by a constant multiple of $x^{-\mathfrak p}(1+x)^{-\mathfrak p}$, integrable on $(0,\infty)$ for $\frac12<\mathfrak p<1$. Dominated convergence and $\int_0^\infty x^{-\mathfrak p}(1+x)^{-\mathfrak p}\mathrm dx=\mathrm B(1-\mathfrak p,2\mathfrak p-1)$ give the claim, and $1-2\mathfrak p>-1$ makes the autocovariance non-summable.
\end{proof}

\begin{proof}[Proof of \cref{thm:env-hayek-horizon}]
(a) $\psi_i^\top\eta_{t+h}=((F^h)^\top\psi_i)^\top\eta_t+\mathfrak f_h$, with $\mathfrak f_h$ a combination of future innovations, independent of $(\eta_t,v_i)$. Hence $\Var(\psi_i^\top\eta_{t+h}\mid s_i,M\eta_t)=\Var(((F^h)^\top\psi_i)^\top\eta_t\mid s_i,M\eta_t)+\Var(\mathfrak f_h)$, and the first term vanishes if and only if $(F^h)^\top\psi_i\in\operatorname{row}(M)$ by \cref{prop:abs-hayek}. The dimension of a linear image of a fixed subspace does not increase under further maps and is preserved by invertible ones, and $(F^{h+1})^\top=F^\top(F^h)^\top$.

(b) Ignoring private signals, the total excess loss of a noiseless message $M\eta_t$ is $\sum_i\Var(\psi_i^\top F^h\eta_t\mid M\eta_t)$, which equals $\tr(\mathsf D_h(I-\mathbb P_M))$ by \cref{lem:env-information}(iii) with a noiseless sensor, $\mathbb P_M$ being the orthogonal projector onto the row space of $MP_0^{1/2}$; and \cref{lem:env-eigen}(b) bounds $\tr(\mathsf D_h\mathbb P_M)$ by $\sum_{\mathsf j\le r_M}\lambda_{\mathsf j}(\mathsf D_h)$. Under detailed balance $P_0^{1/2}(F^h)^\top=\mathsf F^hP_0^{1/2}$, so $\mathsf e_{\mathsf j}^\top\mathsf D_h\mathsf e_{\mathsf j}=\lambda_{\mathsf j}(F)^{2h}\mathsf w_{\mathsf j}$, $\tr\mathsf D_h=\sum_{\mathsf j}\lambda_{\mathsf j}(F)^{2h}\mathsf w_{\mathsf j}$, and the message reading $\mathsf e_1^\top P_0^{-1/2}\eta_t,\dots,\mathsf e_{r_M}^\top P_0^{-1/2}\eta_t$ leaves $\sum_{\mathsf j>r_M}\lambda_{\mathsf j}(F)^{2h}\mathsf w_{\mathsf j}$. Bounding the numerator by $\lvert\lambda_{r_M+1}(F)\rvert^{2h}\sum_{\mathsf j>r_M}\mathsf w_{\mathsf j}$ and the denominator below by $\lvert\lambda_1(F)\rvert^{2h}\mathsf w_1$, which is positive because $\lambda_1(F)\ne0$ and $\mathsf w_1>0$, gives the ratio bound. For the count, whose threshold is defined because $h\ge1$, let $r_M$ be the number of modes at or above the threshold; every other mode has $\lvert\lambda_{\mathsf j}(F)\rvert^{2h}<\lvert\lambda_1(F)\rvert^{2h}\mathfrak e\mathsf w_1/\sum_{\mathsf j'}\mathsf w_{\mathsf j'}$, so the excess is below $\mathfrak e\lvert\lambda_1(F)\rvert^{2h}\mathsf w_1\le\mathfrak e\tr\mathsf D_h$. As $h\to\infty$ the threshold tends to $\lvert\lambda_1(F)\rvert$.

(c) $F^h=\sum_{\mathsf j}\lambda_{\mathsf j}(F)^h\mathsf r_{\mathsf j}\mathsf u_{\mathsf j}^\top$, so $\psi^\top F^h\eta_t=\sum_{\mathsf j}\lambda_{\mathsf j}(F)^h(\psi^\top\mathsf r_{\mathsf j})\mathsf u_{\mathsf j}^\top\eta_t$. Given the slow amplitudes, which determine the terms with $\mathsf j\le r_M$ (conjugate pairs are not split because $\lvert\lambda_{r_M}(F)\rvert>\lvert\lambda_{r_M+1}(F)\rvert$), the conditional variance is at most the variance of the remainder $\sum_{\mathsf j>r_M}\lambda_{\mathsf j}(F)^h(\psi^\top\mathsf r_{\mathsf j})\mathsf u_{\mathsf j}^\top\eta_t$, whose standard deviation is at most the sum of the terms' root-mean-square values, $\lvert\lambda_{\mathsf j}(F)\rvert^h\lvert\psi^\top\mathsf r_{\mathsf j}\rvert(\mathsf u_{\mathsf j}^*P_0\mathsf u_{\mathsf j})^{1/2}$. When $\lambda_1(F)$ is real, simple and strictly dominant, $\Var(\psi^\top F^h\eta_t)=\lambda_1(F)^{2h}(\psi^\top\mathsf r_1)^2\mathsf u_1^\top P_0\mathsf u_1(1+o(1))$.
\end{proof}

\begin{proof}[Derivations for \cref{ex:env-firm}]
With $\mathsf f_1=F_1^h$ and $\mathsf f_2=F_2^h$, the decision variable is $\mathsf f_1\eta_1+\mathsf f_2\eta_2=\frac{\mathsf f_1+\mathsf f_2}2(\eta_1+\eta_2)+\frac{\mathsf f_1-\mathsf f_2}2(\eta_1-\eta_2)$, and $\Var(\eta_1-\eta_2\mid\eta_1+\eta_2)=4P_{0,1}P_{0,2}/(P_{0,1}+P_{0,2})$ gives $\mathcal L_{\mathrm{spot}}$; $\mathcal L_{\mathrm{slow}}$ is the variance of the unread term $\mathsf f_2\eta_2$, and the futures price reads the decision variable itself. The ratio $\mathcal L_{\mathrm{spot}}/\mathcal L_{\mathrm{slow}}=(\mathsf f_1/\mathsf f_2-1)^2P_{0,1}/(P_{0,1}+P_{0,2})$ increases with $h$ and equals one where $\mathsf f_1/\mathsf f_2=(F_1/F_2)^h=1+(1+P_{0,2}/P_{0,1})^{1/2}$, which is $h^\times$. With price noise $\mathsf n^{\mathrm p}_t$ the spot loss is $\mathsf f_1^2P_{0,1}+\mathsf f_2^2P_{0,2}-(\mathsf f_1P_{0,1}+\mathsf f_2P_{0,2})^2/(P_{0,1}+P_{0,2}+\Var(\mathsf n^{\mathrm p}_t))$. The optimal output and the excess loss follow from $\E[p_{t+h}\mid\eta_t]=\mathsf f_1\eta_{1,t}+\mathsf f_2\eta_{2,t}$: the expected profit at output $\E[p_{t+h}\mid\text{readings}]/\mathsf k_a$ is $\E[\E[p_{t+h}\mid\text{readings}]^2]/(2\mathsf k_a)$, and the readings are functions of $\eta_t$ plus independent noise.
\end{proof}

\begin{proof}[Proof of \cref{lem:env-composition}]
Since $\E[\mathsf X_1\mid\mathsf X_2,\mathsf X_3]=\E[\mathsf X_1\mid\mathsf X_2]=\Sigma_{12}\Sigma_{22}^{-1}\mathsf X_2$, $\Sigma_{13}=\E[\E[\mathsf X_1\mid\mathsf X_2]\mathsf X_3^\top]=\Sigma_{12}\Sigma_{22}^{-1}\Sigma_{23}$, and whitening gives $\mathsf C_{13}=\mathsf C_{12}\mathsf C_{23}$. Then $\mathsf C_{13}\mathsf C_{13}^\top=\mathsf C_{12}\mathsf C_{23}\mathsf C_{23}^\top\mathsf C_{12}^\top\preceq\sigma_1(\mathsf C_{23})^2\mathsf C_{12}\mathsf C_{12}^\top$, and eigenvalues are monotone in the Loewner order \citep{horn2012matrix}. The rank bound is that of a product.
\end{proof}

\begin{proof}[Proof of \cref{prop:env-multiplex}]
(a) is the chain rule, and for jointly Gaussian variables $I(\mathsf Y;s\mid y)$ is $\frac12\ln\det\Cov(\mathsf Y\mid y)-\frac12\ln\det\Cov(\mathsf Y\mid s,y)$. (b) Gaussian conditional independence holds exactly when $\Cov(\mathsf Y,s^{\mathrm{ph}}\mid y^{\mathrm{agg}})=0$, and this covariance is $\mathsf Z_h^\top\Cov(\eta_t\mid y^{\mathrm{agg}})(H^{\mathrm{ph}})^\top$ because the future innovations in $\mathsf Y$ are independent of $(\eta_t,s^{\mathrm{ph}},y^{\mathrm{agg}})$. With noiseless aggregates $\Cov(\eta_t\mid T^{\mathrm{agg}}\eta_t)(T^{\mathrm{agg}})^\top=0$, so columns of $\mathsf Z_h$ in the row space of $T^{\mathrm{agg}}$ are annihilated. (c) Write $\mathsf Y=\mathsf z^\top\eta_t+\mathfrak f_h$ with $\mathsf z=(F^h)^\top\psi$ and $\mathfrak f_h$ independent. The squared partial correlation of $\mathsf Y$ with $s^{\mathrm{ph}}$ given $y=y^{\mathrm{agg}}$ is $\Cov(\mathsf Y,s\mid y)\Cov(s\mid y)^{-1}\Cov(s,\mathsf Y\mid y)/\Var(\mathsf Y\mid y)$; its numerator equals that for $\mathsf z^\top\eta_t$, which is at most $\Var(\mathsf z^\top\eta_t\mid y)$, and $I(\mathsf Y;s\mid y)$ is minus one half the logarithm of one minus it. In the reversible case with $y$ reading $\mathsf e_1^\top P_0^{-1/2}\eta_t,\dots,\mathsf e_{r_{\mathrm{agg}}}^\top P_0^{-1/2}\eta_t$, the identity $F^h=P_0^{1/2}\mathsf F^hP_0^{-1/2}$ gives $\Var(\mathsf z^\top\eta_t\mid y)=\sum_{\mathsf j>r_{\mathrm{agg}}}\lambda_{\mathsf j}(F)^{2h}(\mathsf e_{\mathsf j}^\top P_0^{1/2}\psi)^2$, and $\Var(\mathsf Y\mid y)\ge\Var(\mathsf Y\mid\eta_t)=\psi^\top(P_0-F^hP_0(F^h)^\top)\psi=\psi^\top P_0^{1/2}(I-\mathsf F^{2h})P_0^{1/2}\psi$, which is positive for $h\ge1$ and $\psi\ne0$ because the symmetric $\mathsf F$ has spectral radius below one; at $h=0$ it vanishes and the displayed bound does not apply. (d) Under the hypotheses $P_0$ and $F^h$ are block diagonal and each $(F^{h_{\mathfrak d}})^\top\psi_{\mathfrak d}$ is supported on one block, so $\mathsf J$ is block diagonal; a block-diagonal symmetric matrix whose blocks have no common eigenvalue has block-pure eigenvectors. Reverse water-filling (\cref{thm:env-canonical}(c)) gives channel $\mathsf j$ the capacity $\frac12\ln^+(\lambda_{\mathsf j}(\mathsf J)/\mathsf v')$; the capacity used increases in every $\lambda_{\mathsf j}(\mathsf J)$ and decreases in $\mathsf v'$, so raising some eigenvalues at a fixed budget raises $\mathsf v'$ and lowers or keeps the capacity of every other channel.
\end{proof}
